% ============================
% Appendix C. Model Estimation Details
% ============================

\section*{Appendix C Model Estimation Details}
\label{app:C}
\doublespacing

This appendix provides additional detail on the estimation of the mixed logit models and the
derivation of exponential and hyperbolic discount parameters used in the main analysis.

\subsection*{C1. Utility Specification}

We adopt a random–utility framework in which individual $i$ evaluates vaccine alternative $j$
in choice task $t$ according to:

\begin{align}
U_{ijt} &= \beta_0
+ \beta_1\,\text{VaccineDelay}_{\text{std},ijt}
+ \beta_2\,\text{VaccineOrigin}_{\text{std},ijt}
+ \beta_3\,\text{VaccineEfficacy}_{\text{std},ijt} \notag\\
&\quad + \beta_4\,\text{SideEffects}_{\text{std},ijt}
+ \beta_5\,\text{CashIncentives}_{\text{std},ijt}
+ \varepsilon_{ijt}, \tag{C1}
\end{align}

where $\varepsilon_{ijt}$ follows an i.i.d.\ Type I extreme-value distribution.  
For comparison, both conditional logit (MNL) and mixed logit (MXL) models were estimated,
but the MXL is the preferred specification because it allows random coefficients to capture
unobserved preference heterogeneity.

In vector notation:

\begin{equation}
U_{ijt}
= \beta_0
+ \sum_{k=1}^{K} \beta_k X_{k,ijt}
+ \varepsilon_{ijt},
\tag{C2}
\end{equation}

Standardization of the attribute coefficients facilitates comparison across attributes and inputs
directly into the discount-rate estimation procedure described below.

\subsection*{C2. Discount Function Specification}

Observed delay levels $t \in \{0,1,3,6\}$ months are mapped to standard exponential and
hyperbolic discounting frameworks:

\begin{align}
D_{\text{Exp}}(t;\hat{\delta}) &= e^{-\hat{\delta} t}, \qquad \hat{\delta} \ge 0, \tag{C3}\\[4pt]
D_{\text{Hyper}}(t;\hat{\kappa}) &= \frac{1}{1+\hat{\kappa} t}, \qquad \hat{\kappa} \ge 0. \tag{C4}
\end{align}


These formulations correspond to time-consistent (exponential) and present-biased
(hyperbolic) discounting, respectively.

\subsection*{C3. Constructing Observed Discount Points}

Let $\beta_{\text{delay (per month)}}$ be the unstandardized coefficient from the secondary model in
which delay is entered linearly. This coefficient represents the marginal utility loss associated
with a one-month increase in delay. It is recovered from the standardized coefficient via:

\[
\beta_{\text{delay (per month)}} 
= \frac{\beta_{\text{delay (std)}}}{\mathrm{SD}(\text{MultipleDelays})}. \tag{C5}
\]

Following standard practice in preventive-health time-preference studies, we define the
utility-consistent “observed” discount factor for each delay $t$ as:

\[
D_{\text{obs}}(t) 
= \exp\!\left(-\beta_{\text{delay (per month)}}\, t\right). \tag{C6}
\]

These points serve as empirical inputs for fitting exponential and hyperbolic discount functions.

\subsection*{C4. Estimation of Discount Parameters}

The exponential and hyperbolic discount parameters are obtained by minimizing the sum of
squared errors (SSE) between the observed and theoretical discount functions.

\paragraph{Hyperbolic fit.}

\[
\mathrm{SSE}_{\text{Hyper}}(\hat{\kappa})
= \sum_{i=1}^{n}
\left[
D_{\text{obs}}(t_i)
- D_{\text{Hyper}}(t_i;\hat{\kappa})
\right]^2 .
\tag{C7}
\]


The first-order condition is:

\[
\frac{\partial}{\partial \kappa}
\mathrm{SSE}_{\text{Hyper}}(\kappa)
= 2 \sum_{i=1}^{n}
\left(
D_{\text{obs}}(t_i)
- \frac{1}{1 + \kappa t_i}
\right)
\frac{t_i}{(1 + \kappa t_i)^2}
= 0,
\qquad \text{evaluated at } \kappa = \hat{\kappa}.
\tag{C8}
\]


\paragraph{Exponential fit}

\[
\mathrm{SSE}_{\text{Exp}}(\hat{\delta})
= \sum_{i=1}^{n}
\left[
D_{\text{obs}}(t_i)
- D_{\text{Exp}}(t_i;\hat{\delta})
\right]^2 .
\tag{C9}
\]

First-order condition:

\[
\frac{\partial}{\partial \delta}
\mathrm{SSE}_{\text{Exp}}(\delta)
= 2 \sum_{i=1}^{n}
\left(
D_{\text{obs}}(t_i)
- e^{-\delta t_i}
\right)
t_i e^{-\delta t_i}
= 0,
\qquad \text{evaluated at } \delta = \hat{\delta}.
\tag{C10}
\]

Closed-form solutions do not exist; parameters were estimated via nonlinear least squares.

\subsection*{C5. Subgroup Estimation and Multiplicity}

Subgroup-specific MXL models were estimated separately for demographic, behavioural,
health-status, and institutional-trust categories. Because this involves multiple hypothesis tests,
$p$-values for subgroup delay coefficients were adjusted using the Benjamini–Hochberg
false-discovery rate (FDR) at $q = 0.10$. Robustness of subgroup differences was confirmed
under Holm–Bonferroni corrections.


\subsection*{C.6 Recovering Discount Parameters from Waiting-Time Coefficients}

To recover exponential discount parameters from the DCE model, we map the estimated 
waiting-time coefficient $\beta_{\text{delay}}$ to the implied monthly discount factor $\delta$. 
Under exponential discounting, utility declines with delay according to $U(t) = U_0 \delta^{t}$. A first-order approximation yields a linear representation in which 
the marginal disutility of delay corresponds to the logarithm of the discount factor. 
Normalizing the baseline utility $U_0 = 1$ implies:

\[
\hat{\delta} = \exp\!\left(\hat{\beta}_{\text{delay}}\right).
\]


Thus, more negative values of $\beta_{\text{delay}}$ correspond to lower discount factors and 
stronger impatience toward delayed protection. This transformation is used to obtain the 
discount parameters reported in Table~\ref{tab:subgroups_long}.

\subsection*{C.7 MWTA Derivation}

The marginal willingness-to-accept (MWTA) for one month of delay is derived as the 
ratio of the delay coefficient to the cash-incentive coefficient:

\[
\text{MWTA} = -\frac{\beta_{\text{delay}}}{\beta_{\text{cash}}}.
\]

This represents the monetary compensation required to offset the disutility of an 
additional month of waiting.



\subsection*{C8. Additional Diagnostics}

We conducted additional diagnostic checks, including model comparison (AIC and
log-likelihood), scale-heterogeneity assessment using a GMNL specification, and
re-estimation excluding potential speeders. Across all alternative specifications, the sign,
magnitude, and statistical significance of the delay coefficient remained highly stable,
and the implied discount parameters were unchanged. These results confirm that the
main findings are robust to functional-form assumptions and response-quality concerns.
Full model outputs are available upon request.



\subsection*{C9. Behavioral Delay Multiplier (BDM)}

To provide an intuitive interpretation of the hyperbolic discount parameter, we compute a
Behavioral Delay Multiplier (BDM), defined as:

\[
\mathrm{BDM} = \frac{1}{1 - \hat{\kappa}}, \tag{C11}
\]

which rescales the estimated hyperbolic parameter into a measure of the perceived cost of
waiting. A value greater than 1 indicates that each month of objective delay is experienced as
a longer \emph{behavioural} delay. For example, the overall estimate $\hat{\kappa} = 0.225$
implies a BDM of 1.29, meaning that one month of waiting feels equivalent to approximately
1.29 months. Subgroups with larger $\hat{\kappa}$—such as older adults, respondents with
primary education, or those reporting neutral trust in government—exhibit BDM values
between 1.45 and 1.55, indicating substantially amplified perceived waiting costs. In contrast,
groups with smaller $\hat{\kappa}$, including rural residents and individuals reporting regular
physical activity, show BDM values near 1.15–1.20, reflecting weaker behavioural sensitivity
to delay. This multiplier offers a transparent summary measure that complements the formal
discount-rate estimates presented in Sections~C2–C4.

\begin{table}[H]
\centering
\caption{Behavioral Delay Multiplier (BDM) for Selected Subgroups}
\label{tab:BDM}
\begin{tabular}{lccc}
\toprule
\textbf{Subgroup} & $\hat{\kappa}$ & \textbf{BDM} & \textbf{Interpretation} \\
\midrule
Overall sample          & 0.225 & 1.29 & 1 month feels like 1.29 months \\
Urban residents         & 0.290 & 1.41 & Higher perceived waiting cost \\
Rural residents         & 0.174 & 1.21 & Lower perceived waiting cost \\
Primary education       & 0.348 & 1.53 & Amplified behavioural delay \\
Adults aged 65+         & 0.327 & 1.49 & Delay feels substantially longer \\
Neutral trust in gov.   & 0.356 & 1.55 & Strongest impatience \\
\bottomrule
\end{tabular}
\end{table}


